\chapter{总结与展望}
\label{ch:22}

\paragraph*{本章阅读锚点}
本章的锚点是技术路线回看。读这一章时不要再记单篇方法，而要把全书方法归入表示、优化、数据、规模、应用和安全几个长期问题。


如前言所述，3DGS 从一项渲染技术膨胀为一个横跨表示学习、几何重建、材质分解、动态建模、场景生成、语义理解、模型压缩、大规模建图、机器人导航、自动驾驶、医学影像、遥感测绘、安全隐私等领域的完整研究生态系统。

本书系统梳理了这一生态系统的 21 个核心方向，覆盖 2679 篇论文。本章从三个角度为全书收尾。第~\ref{sec:22_1}~节回顾全书的六大部分和核心发现，建立一张从表示到应用的技术全景图；第~\ref{sec:22_2}~节提炼横跨多个方向的共性技术趋势，揭示 3DGS 研究正在经历的深层范式转变；第~\ref{sec:22_3}~节面向未来，讨论十个尚未解决的开放问题。

\begin{figure}[htbp]
\centering
\resizebox{\textwidth}{!}{\begin{tikzpicture}[
  route/.style={rectangle, rounded corners=5pt, draw=#1, fill=#1!10,
    text width=3.2cm, minimum height=0.9cm, align=center, font=\small},
  method/.style={rectangle, rounded corners=3pt, draw=#1, fill=#1!6,
    font=\scriptsize, text width=3.0cm, align=center, minimum height=0.6cm},
  arr/.style={-{Stealth[length=4pt]}, techgray!60},
  >=Stealth
]

  \node[route=techgray, text width=4cm, font=\small\bfseries] (root) at (12.0, 5.6) {3DGS 技术全景与趋势};
  \node[route=techblue] (b0) at (5.8, 3.1) {研究全景回顾};
  \node[method=techblue] (b0_l0) at (1.7, 0.6) {表示与渲染\\{\tiny Mip-Splatting, 2DGS, 3DGS-MCMC}};
  \node[method=techblue] (b0_l1) at (5.8, 0.6) {效率与泛化\\{\tiny Scaffold-GS, LightGaussian, DUSt3R}};
  \node[route=techgreen] (b1) at (12.0, 3.1) {技术趋势};
  \node[method=techgreen] (b1_l0) at (9.9, 0.6) {大模型融合\\{\tiny LRM, GS-LRM, VGGT}};
  \node[method=techgreen] (b1_l1) at (14.0, 0.6) {规模与建图\\{\tiny CityGaussian, SplaTAM}};
  \node[route=techpurple] (b2) at (17.9, 3.1) {开放问题};
  \node[method=techpurple] (b2_l0) at (17.9, 0.6) {长尾场景\\{\tiny 稀疏视角, 极端光照}};
  \node[method=techpurple] (b2_l1) at (21.5, 0.6) {可信可靠\\{\tiny 水印鲁棒性, 隐私保护}};
  \draw[arr] (root) -- (b0);
  \draw[arr] (root) -- (b1);
  \draw[arr] (root) -- (b2);
  \draw[arr] (b0) -- (b0_l0);
  \draw[arr] (b0) -- (b0_l1);
  \draw[arr] (b1) -- (b1_l0);
  \draw[arr] (b1) -- (b1_l1);
  \draw[arr] (b2) -- (b2_l0);
  \draw[arr] (b2) -- (b2_l1);
\end{tikzpicture}}
\caption{3DGS 技术全景与趋势}
\label{fig:ch22_outlook_taxonomy}
\end{figure}

%=====================================================================
\section{研究全景回顾}
\label{sec:22_1}

全书按六部分组织，遵循"表示与渲染 $\to$ 几何与物理 $\to$ 生成与理解 $\to$ 效率与泛化 $\to$ 规模与建图 $\to$ 应用领域"的技术演进逻辑。

\paragraph{Part I：表示与渲染——3DGS 的技术根基}
第一部分（第~\ref{ch:01}--\ref{ch:03}~章）建立了全书的技术基础。

第~\ref{ch:01}~章从 NeRF 的体渲染瓶颈出发，详细拆解了 3DGS 的完整技术栈：高斯基元的参数化定义、3D-to-2D 投影与 $\alpha$-混合渲染、自适应密度控制（ADC）机制，以及基于 tile 的可微光栅化器。该章特别强调了 NeRF 体渲染与 3DGS $\alpha$-混合在数学上的等价性——这一等价性解释了 3DGS 为何能在保持渲染质量的同时获得三个数量级的速度提升。此外，该章还梳理了领域综述的分类学框架、基准数据集、理论分析和开源工具链。

第~\ref{ch:02}~章是全书论文最密集的章节（286 篇），从渲染质量角度全面优化 3DGS 管线。抗锯齿（Mip-Splatting）、密度控制重构（3DGS-MCMC）、训练加速（3DGS-LM）、光线追踪（3DGRT、EVER）、硬件加速（Lumina ASIC）等改进，将 3DGS 的渲染质量推向了 NeRF 时代最佳方法（Mip-NeRF 360）的水平。同时，训练时间从小时级压缩到分钟级。多模态渲染（事件相机、热红外、LiDAR、音频）和超分辨率技术进一步拓展了 3DGS 的适用边界。

第~\ref{ch:03}~章聚焦高斯基元本身的表示设计，从标准 3D 椭球出发，发展出 2D 高斯（2DGS）、广义高斯（Gaussian Shading）、神经高斯（Neural Gaussians）等变体，每种变体针对特定的几何或渲染需求进行了表示层面的优化。

\paragraph{Part II：几何、材质与动态——从外观到物理世界}
第二部分（第~\ref{ch:04}--\ref{ch:06}~章）将 3DGS 从"好看"推向"准确"和"动态"。

第~\ref{ch:04}~章解决的核心问题是：高斯基元是体积性的，没有表面概念，如何从中提取精确几何？五条技术路线——深度与法线正则化、SDF（Signed Distance Function，符号距离函数）引导、多视角几何约束、混合 Mesh-Gaussian 表示、网格提取——逐步从"正则化引导的高斯优化"演进到"几何与渲染的统一表示"。最新工作（MILo、OMeGa、Gaussian-Voxel Duet）不再将几何和渲染视为两个独立任务，而是在统一框架中联合优化。

第~\ref{ch:05}~章将 3DGS 从 Lambertian 漫反射推向复杂材质和逆渲染。球谐函数无法表达的强视角依赖性（镜面反射、透明折射、次表面散射），通过 PBR（Physically-Based Rendering，基于物理的渲染）参数化、可微光线追踪和分解式表示得到系统解决。重光照、材质编辑和场景级逆渲染是当前的前沿方向。

第~\ref{ch:06}~章是全书第二大章（250 篇），将 3DGS 从静态推向动态。时间建模从简单的全局变形场，发展为规范空间加变形场的解耦架构、4D 高斯表示（4DGS）、运动学骨架驱动、神经场变形等多元路线。动态场景重建已成为 3DGS 最活跃的研究方向之一。

\paragraph{Part III：生成、编辑与理解——从重建到创造}
第三部分（第~\ref{ch:07}--\ref{ch:10}~章）将 3DGS 的角色从"世界观测工具"扩展为"内容创造工具"和"场景理解工具"。

第~\ref{ch:07}~章将物理仿真引入 3DGS，从紧耦合（MPM 物质点法）到松耦合（代理结构），再到物理属性感知和声音合成，使 3DGS 场景不仅可渲染，还可交互、可仿真。物理精度的验证标准、多材料耦合和实时交互仍是核心挑战。

第~\ref{ch:08}~章覆盖了 3DGS 在 3D 生成领域的完整技术栈。Score 蒸馏（SDS，Score Distillation Sampling，及其变体）将 2D 扩散先验蒸馏到 3DGS 表示中；前馈重建模型（LRM、VGGT）实现秒级生成；3D 生成模型（GAN、扩散、自回归、VAE）直接在 3DGS 空间做采样生成，正在快速逼近 2D 生成模型的质量和多样性。

第~\ref{ch:09}~章展示了 3DGS 场景如何被编辑和操控——文本引导的语义编辑、交互式物体操控、风格化、场景修复和几何变形。编辑精度与多视角一致性的矛盾是核心挑战。

第~\ref{ch:10}~章将高斯基元从颜色载体扩展为语义、语言和实例信息的载体，实现了从"看得见"到"看得懂"的跃迁。2D 基础模型（SAM、CLIP、DINOv2）的语义能力通过特征蒸馏提升到 3D 空间。开放词汇分割、3D 目标检测和场景图构建使 3DGS 场景具备了语义理解能力。

\paragraph{Part IV：效率与泛化——从实验室到产品}
第四部分（第~\ref{ch:11}--\ref{ch:13}~章）解决 3DGS 走出实验室的核心瓶颈。

第~\ref{ch:11}~章从压缩与加速的角度系统优化 3DGS 的存储和计算效率。剪枝（LightGaussian）、量化与熵编码（HAC++、ContextGS）、锚点表示（Scaffold-GS）、流式渲染（Matryoshka GS）等技术的组合，已将 3DGS 模型从数 GB 压缩到数 MB，同时保持可接受的渲染质量。从"训练后压缩"到"训练中控制"，从"固定码率"到"可伸缩码率"，从"逐场景优化"到"通用前馈"——这些趋势共同推动 3DGS 走向手机和 VR 头显等边缘设备。

第~\ref{ch:12}~章解决 3DGS 在非理想输入条件下的鲁棒性问题——稀疏视角、未知位姿、图像退化、瞬态干扰物。基础模型（DUSt3R、MASt3R）和生成式先验（扩散模型）正从根本上改变稀疏视角重建的范式，从被动优化转向利用强先验的主动重建。此外，主动视觉进一步将 3DGS 从静态表示转变为感知-决策闭环系统。

第~\ref{ch:13}~章探索可泛化前馈重建——训练一个通用模型，推理时直接从图像预测高斯参数。从像素对齐到体素对齐、从固定密度到自适应分配、从单次推断到迭代精炼，前馈方法正在缩小与逐场景优化方法之间的质量差距。无标定方法（NoPoSplat、Splatt3R）和大重建模型（VGGT-X 可处理 1000+ 视角）代表了该方向的前沿。

\paragraph{Part V：规模与建图——从物体到城市}
第五部分（第~\ref{ch:14}--\ref{ch:15}~章）将 3DGS 的规模从室内场景推到城市级别。

第~\ref{ch:14}~章处理大规模场景的分治策略、在线增量重建、层次化 LOD（Level of Detail，细节层次）表示和分布式并行训练。CityGaussian、BlockGaussian 等方法将分治策略标准化；TideGS 在单 GPU 上训练十亿级高斯；虚拟内存技术使消费级硬件能够处理城市级场景。从离线优化到在线服务，从单机到分布式——大规模 3DGS 正在成为城市数字孪生的技术基础。

第~\ref{ch:15}~章将 3DGS 引入 SLAM（Simultaneous Localization and Mapping，同步定位与建图）与稠密建图，带来范式转换：地图从稀疏点云变为可实时渲染的照片级三维场景。3DGS-SLAM 在跟踪精度和重建质量上已展现出明显优势。从 RGB-D 到单目 RGB，从纯视觉到多模态融合，从静态到动态——技术路线日趋成熟。长期运行与终身建图、边缘部署和语义-几何联合优化是下一步的关键挑战。

\paragraph{Part VI：应用领域——从通用到专用}
第六部分（第~\ref{ch:16}--\ref{ch:21}~章）展示了 3DGS 在六个垂直领域的深度应用。

第~\ref{ch:16}~章覆盖人体化身与动画，从人脸重建到全身化身、从说话人头像到物理可信的衣物动画，3DGS 的实时渲染和显式表示为数字人技术提供了理想基础。第~\ref{ch:17}~章展示 3DGS 在机器人导航与具身智能中的变革性应用：3DGS 仿真生成的训练数据使机器人策略实现零样本 Sim-to-Real 迁移（SplatSim 达到 86.25\% 成功率），闭环具身智能和多模态统一表示是核心挑战。第~\ref{ch:18}~章聚焦自动驾驶场景，从前馈 4D 街景重建到闭环仿真平台，3DGS 正在成为自动驾驶数据引擎的核心技术。

第~\ref{ch:19}~章将 3DGS 带入医学影像，从内窥镜到 CT/MRI 重建、从手术导航到显微成像，3DGS 的实时渲染和紧凑表示在多个医学子领域展现出独特价值。第~\ref{ch:20}~章覆盖遥感与航拍，EOGS 等框架将 3DGS 适配到卫星 RPC 相机模型，实现从天级到分钟级的重建加速。第~\ref{ch:21}~章讨论 3DGS 的安全与隐私问题，包括数字水印、对抗攻防和信息隐藏，揭示了 3DGS 在知识产权保护方面的脆弱性和防御策略。

\begin{takeaway}
研究全景回顾覆盖全书 2679 篇论文的六大技术板块：从表示与渲染的基础理论，到动态场景与几何处理，再到生成、理解与编辑，最后深入六个垂直领域应用，构成了从通用到专用的完整技术地图。
\end{takeaway}

%=====================================================================
\section{技术趋势}
\label{sec:22_2}

回顾全书 2679 篇论文，若干共性趋势横跨多个章节，揭示了 3DGS 研究正在经历的深层范式转变。

\paragraph{从隐式到显式再到混合表示}
3DGS 的核心贡献是用显式高斯基元替代 NeRF 的隐式 MLP。然而，三年的研究表明，纯粹的显式表示并非银弹。第~\ref{ch:04}~章中 SDF 引导的表面重建、第~\ref{ch:05}~章中神经材质场的引入、第~\ref{ch:11}~章中锚点加神经场的混合表示（Scaffold-GS），都指向同一个趋势：未来的 3DGS 系统将是显式高斯与隐式场的混合体。显式高斯负责高效渲染和可编辑性，隐式场负责表达高维连续函数（材质、语义、变形场）。如何在统一框架中优雅地融合两者，是表示层面的核心问题。

\paragraph{从逐场景优化到前馈泛化}
全书最显著的趋势之一是"前馈化"。第~\ref{ch:13}~章的可泛化前馈重建、第~\ref{ch:08}~章的 3D 生成模型、第~\ref{ch:11}~章的通用压缩编解码器，都在追求同一个目标：一个训练好的通用模型在新场景上直接推理，无需逐场景优化。DUSt3R/MASt3R 等几何基础模型、VGGT/LRM 等重建模型、Stable Diffusion 等生成模型的引入，为这一趋势提供了技术基础。当前前馈方法在渲染质量上仍落后于逐场景优化，但差距正在快速缩小。

\paragraph{从重建工具到感知基础架构}
3DGS 最初是为新视角合成而设计的。但第~\ref{ch:10}~章的场景理解、第~\ref{ch:15}~章的 SLAM、第~\ref{ch:17}~章的具身智能、第~\ref{ch:18}~章的自动驾驶感知，都表明 3DGS 正在从"重建工具"演进为"感知基础架构"——一种既可渲染又可查询、既有几何又有语义的统一场景表示。高斯基元从纯粹的颜色载体扩展为语义、语言、实例和物理属性的载体。这使得 3DGS 场景成为下游任务（导航、操控、检测、分割）的直接输入。

\paragraph{从静态到动态到交互}
时间维度的引入是 3DGS 近两年最活跃的方向之一。第~\ref{ch:06}~章的 4DGS 动态建模、第~\ref{ch:07}~章的物理仿真、第~\ref{ch:16}~章的人体动画，将 3DGS 从"冻结的三维快照"推向"可交互的动态世界"。但动态建模的复杂度远超静态：时序一致性、运动物体的分解与追踪、物理约束的引入、4D 数据的存储和流式传输，每一个子问题都是独立的研究方向。实时交互的物理仿真以及 4D 场景的编辑和生成，仍处于早期阶段。

\paragraph{从实验室到产品化}
效率优化（第~\ref{ch:11}~章）和大规模部署（第~\ref{ch:14}~章）的趋势表明，3DGS 正在从学术概念验证走向工程产品化。专用硬件加速器（Lumina ASIC、3DGauCIM）、流式渲染协议、压缩编解码标准、跨平台工具链——这些基础设施的成熟，将决定 3DGS 能否在手机、VR/AR 头显、汽车和云端实现大规模部署。从研究论文数量趋势看，效率与部署相关的工作增速相对较快。这反映了社区的优先级正在从"做得更好"转向"做得更快更小更实用"。

\paragraph{基础模型驱动的范式转移}
DUSt3R/MASt3R（几何）、SAM（分割）、CLIP/DINOv2（语义）、Stable Diffusion（生成）等 2D/3D 基础模型的引入，正在从根本上改变 3DGS 研究的范式。基础模型提供强先验，3DGS 提供高效表示——两者的结合使得极稀疏视角重建（第~\ref{ch:12}~章）、开放词汇场景理解（第~\ref{ch:10}~章）、文本到 3D 生成（第~\ref{ch:08}~章）等此前被认为极度困难的问题变得可解。未来，3DGS 很可能不再作为一个独立的技术栈存在，而是作为基础模型生态系统中的一环——一种高效的 3D 表示层，由上下游的基础模型所驱动。

\begin{takeaway}
技术趋势揭示了 3DGS 正在经历的深层范式转变：从纯显式到显式-隐式混合表示、从逐场景优化到前馈泛化、从重建工具到感知基础架构、从实验室到产品化，基础模型的融合正在重新定义 3DGS 的技术定位。
\end{takeaway}

%=====================================================================
\section{开放问题}
\label{sec:22_3}

尽管 3DGS 在三年内取得了令人瞩目的进展，若干根本性问题仍未解决。以下十个开放问题横跨全书多个章节，代表了未来研究最有价值的方向。

\paragraph{理论基础的薄弱}
第~\ref{ch:01}~章的理论分析揭示了 3DGS 缺乏严格的收敛性保证、泛化界和最优高斯数量的理论指导。ADC 的启发式判据（位置梯度阈值、$\beta$ 参数）至今没有理论上的最优选择。为什么数百万个高斯在优化后能如此好地近似复杂的辐射场？这个看似简单的问题缺乏严格的理论解释。建立 3DGS 的数学基础——类似于变分自编码器（VAE）的 ELBO 界或扩散模型的去噪理论——是该领域走向成熟的前提。

\paragraph{几何精度的根本瓶颈}
第~\ref{ch:04}~章详细讨论了几何与表面重建的挑战。高斯基元的体积性本质——没有"内""外"之分——使得精确表面提取始终是一个近似过程。尽管 SDF 引导和混合 Mesh-Gaussian 表示取得了显著进展，但在细杆结构（thin structures）、透明/半透明表面和复杂拓扑方面，3DGS 的几何精度仍不如传统 MVS（Multi-View Stereo，多视角立体匹配）和 LiDAR 方法。如何在不牺牲渲染效率的前提下获得亚毫米级的几何精度，是连接 3DGS 与工业级应用（如 CAD、数字孪生）的关键桥梁。

\paragraph{动态场景的 4D 一致性}
第~\ref{ch:06}~章中 4DGS 的快速发展掩盖了一个根本问题：如何保证动态重建的时序一致性？当前方法在短序列（几秒到几十秒）上效果良好，但在长时间跨度内累积的漂移、外观变化和运动建模误差会使重建质量显著退化。物理约束的引入（第~\ref{ch:07}~章）部分缓解了这一问题，但以牺牲灵活性为代价。一个真正鲁棒的 4D 表示，需要在表达能力和物理一致性之间找到更好的平衡。

\paragraph{物理仿真的真实性与验证}
第~\ref{ch:07}~章的物理仿真方法在视觉质量上日益逼真，但视觉真实不等于物理真实。大多数方法仅通过 PSNR/SSIM 评估仿真结果，缺少基于物理度量的评估基准。碰撞响应是否正确？材料断裂模式是否符合物理？流体的雷诺数和涡旋结构是否合理？这些问题在没有标准化物理评估基准的情况下无法回答。建立仿真结果的物理验证体系，是 3DGS 物理仿真走向工程应用的先决条件。

\paragraph{生成质量与可控性}
第~\ref{ch:08}~章的 3D 生成技术在多样性上快速进步，但与 2D 生成模型相比仍有显著差距。更深层的问题是可控性：当前方法难以精确控制生成结果的几何结构、材质属性和物理合理性。场景级生成的全局一致性（多物体交互、复杂布局）仍是一个未解难题。此外，3D 训练数据的规模和质量远不及 2D 领域——互联网规模的 3D 资产数据集至今仍不存在。

\paragraph{编辑的精度与一致性}
第~\ref{ch:09}~章揭示了 3DGS 场景编辑的根本矛盾：扩散模型的大可行输出空间使精确编辑难以保持多视角一致。当前的缓解策略（渐进式分解、蒸馏方法）在复杂编辑（如局部形变后保持全局光照一致）上仍有局限。理想的编辑系统应具备可逆性（撤销/重做）、物理一致性（形变后纹理不拉伸、光照匹配）和层级性（物体级、部件级、材质级编辑的统一框架）。

\paragraph{跨场景泛化的语义理解}
第~\ref{ch:10}~章的场景理解方法大多需要逐场景优化，无法直接迁移到新场景。如何学习通用的 3D 语义表示——使得一个训练好的模型能在任何新场景上直接进行开放词汇分割、目标检测和场景图构建——是通向大规模部署的关键。细粒度语义（部件级、材质级、属性级）和动态场景的在线语义理解是两个尚未被充分探索的前沿方向。

\paragraph{大规模实时渲染的系统瓶颈}
第~\ref{ch:14}~章在单 GPU 上训练十亿级高斯已可实现，但实时渲染仍是瓶颈。城市级场景的核外渲染、分布式训练的收敛性保证、持续终身更新（场景在数月甚至数年变化中的模型维护）、异构设备（云端 GPU 到边缘设备）的协作渲染——这些系统工程问题决定了 3DGS 能否成为城市数字孪生的基础设施。

\paragraph{Sim-to-Real 迁移的领域差距}
第~\ref{ch:07}~章的物理仿真、第~\ref{ch:17}~章的机器人仿真和第~\ref{ch:18}~章的自动驾驶仿真都面临同一个问题：3DGS 重建的场景在视觉上逼真，但物理属性（摩擦系数、弹性模量、质量分布）缺失或不准确。在仿真中训练的策略迁移到真实世界时，领域差距导致性能下降。Real2Sim2Real 框架和物理感知的重建是缓解策略。但根本解决方案需要将物理属性作为一等公民集成到 3DGS 表示中。

\paragraph{安全与隐私的制度化}
第~\ref{ch:21}~章揭示了 3DGS 在安全领域的双重身份——既是需要被保护的对象（数字水印、对抗防御），也是可用于攻击的工具（对抗样本生成）。当前水印方法缺乏理论上可证明的安全性，攻防军备竞赛仍在持续。随着 3DGS 资产的商业化（数字人、建筑模型、城市孪生），知识产权保护和隐私合规将成为制度化的需求——但标准化的评估基准和威胁模型定义尚未建立。

\begin{takeaway}
开放问题指向 3DGS 未来最有价值的研究方向：理论基础的夯实、物理精度与 Sim-to-Real 迁移的领域差距、安全隐私的制度化需求——从"能不能用"到"用得足够好、足够快、足够安全"的跨越。
\end{takeaway}

\section*{本章小结}

本书从 3DGS 的技术基础出发，经过 21 个章节的系统梳理，覆盖了 2679 篇论文，建立了一张从表示到渲染、从几何到材质、从静态到动态、从重建到生成、从理解到编辑、从效率到规模、从通用到专用的完整技术地图。3DGS 的三年发展史证明了一件事：一项足够优雅的技术创新（用显式高斯替代隐式场），不仅能解决眼前的工程问题（渲染速度），还能催生一个完整的研究生态系统。

展望未来，3DGS 面临的挑战不再是"能不能用"，而是"能不能用得足够好、足够快、足够安全"。理论基础的夯实、物理精度的提升、基础模型的融合、大规模系统的工程化、安全隐私的制度化——这些问题的逐步解决，将推动 3DGS 从一项渲染技术成长为空间智能的基础设施。
